\scanpage{298}
\margin{299}
\begin{center}- 34 -\end{center}\hfill IX

On observera que $S'$ est noethérien régulier connexe de dimension 1, ce qui donne un sens à la conclusion énoncée. Pour prouver 3.6, choisissons une clôture séparable $\overline K$ de $K$, et écrivons l'extension $\overline K$ comme limite inductive de ses sous-extensions galoisiennes finies $K_i$. Soit $S_i$ le normalisé de $S$ dans $K_i$, et soit $U$ le plus grand ouvert de $S$ sur lequel $A_K$ ait bonne réduction. Comme $S$ est noethérien connexe de dimension 1, et $U\ne\varnothing$, on voit que $T=S-U$ est un ensemble fini de point fermés de $S$, et que pour tout $t\in T$, $\underline O_{S,t}$ est un anneau de valuation discrète. Il suffit de trouver, pour tout $t\in T$, un indice $i=i(t)$ tel que $A_{K_i}$ ait réduction semi-stable en les points de $S_i$ au-dessus de $t$ ; en effet, il suffira de prendre alors pour $K'$ le composé des extensions $K_{i(t)}$ pour $t\in T$. Cela nous ramène donc au cas où $S$ est un trait. D'autre part, comme de groupe $\operatorname{Gal}(K'/K)$ permute entre eux les points de $S_i$ au-dessus du point fermé $s$ de $S$, il suffit de trouver $K'$ de telle façon que $A_{K'}$ ait réduction semi-stable en \uline{un} point du normalisé $S'$ ; d'ailleurs il est inutile ici que $K'$ soit galoisien sur $K$, il suffit de trouver $K'$ extension finie étale de $K$ satisfaisant la condition précédente, car la clôture galoisienne de cette extension satisfera à la même condition. Utilisant le critère 3.5 (iv) on voit donc que 3.6 revient à l'assertion suivante :

\uline{Corollaire} 3.7. \uline{On suppose à nouveau} $S$ \uline{un trait}. \uline{Alors il existe un sous-groupe ouvert} $\pi'$ \uline{de} $\pi$ \uline{tel que l'action} $I'=I\cap\pi'$ \uline{sur} $U=T_\ell(A(\overline K))$ \uline{soit unipotente d'échelon} 2, \uline{ou encore} : \uline{il existe un sous-groupe ouvert} $I'$ \uline{de} $I$ \uline{dont l'action sur} $U$ \uline{soit unipotente d'échelon} 2.

\margin{352}
\margin{352}
